\chapter{客观要件二：结果}

% Source PDF physical page 54; see audit/source_coverage.json.
\begin{EditionOpeningPanel}
1.行为人实施了危害行为，接下来考察行为人有无制造法益侵害事实（违法事实）。法益侵害事实包括两种样态，一是危险状态，二是实害结果。

2.重要考点：结果加重犯的因果关系（“因”的辨认）。
\end{EditionOpeningPanel}

\section{法益侵害事实的分类}

[问题]狗蛋用面粉冒充胃药，进行销售。由于成分是面粉，故对身体没有造成危害。狗蛋是否构成生产、销售假药罪？（假药案，2010年第15题）

\subsection*{一、危险犯与实害犯}

1.实害犯。立法者规定，某个犯罪的成立需要具备实害结果，这种犯罪称为实害犯，也称为结果犯。例如，立法者规定，生产、销售劣药罪的成立要求有实害结果“对人体健康造成严重危害”。该罪便是实害犯。这种实害犯也称为结果犯。

2.具体危险犯。立法者规定，某个犯罪的成立只需要具备具体危险，这种犯罪称为具体危险犯。例如，立法者规定，生产、销售不符合标准的医用器材罪的成立，要求产生“足以严重危害人体健康”的危险，这种危险需要法官具体判断。该罪便是具体危险犯。又如，放火罪、投放危险物质罪、以危险方法危害公共安全罪（第114条），均是具体危险犯。（2015年第14题）

3.抽象危险犯。立法者规定，某个犯罪的成立只需要具备抽象危险，这种犯罪称为抽象危险犯。对于这种犯罪，立法者只描述行为，认为实施了该行为就会产生抽象危险，这种危险是立法者认定的，不需要法官具体判断（2017年第2题）。例如，立法者规定，生产、销售假药的，构成生产、销售假药罪。立法者认为，只要实施了这种行为，就会有一种抽象的危险。该罪属于抽象危险犯。抽象危险犯也称为行为犯。因此，在本节开头“假药案”中，狗蛋的假药虽然没有给人的身体造成实害，也没有具体危险，但有可能贻误病情，推定有抽象危险，因此构成生产、销售假药罪。（2010年第15题）

% Source PDF physical page 55; see audit/source_coverage.json.
\subsection*{二、行为犯与结果犯}

1.行为犯

这是指不将实害结果作为犯罪成立条件的犯罪。例如，生产、销售假药罪，生产、销售有毒、有害食品罪，伪证罪，诬告陷害罪，组织、领导、参加黑社会性质组织罪。

2.结果犯

这是指将实害结果作为犯罪成立条件的犯罪。例如，丢失枪支不报罪，滥用职权罪，生产、销售劣药罪等，这些犯罪要成立，法条规定要求造成实害结果。这些犯罪被称为结果犯，也称为实害犯。（2017年第2题）

\section{结果加重犯}

[问题］《水浒传》里，鲁智深拳打镇关西，只想教训，未料打死。对鲁智深该如何处理？

\InsertDiagram{figures/manual_p055_1.tex}

结果加重犯，是指一个行为构成基本犯，该行为同时导致了加重结果，基于此，刑法对基本犯加重处罚。“故意伤害罪致人死亡”是典型的结果加重犯。例如，甲欲伤害乙，用刀刺伤乙的腿部，造成轻伤，乙因失血过多而死亡。甲的伤害行为构成故意伤害罪。该伤害行为也导致死亡结果，甲对死亡结果持过失心理。甲构成结果加重犯“故意伤害罪致人死亡”。又如，鲁智深拳打镇关西，构成故意伤害罪致人死亡。

\subsection*{一、法定性}

不能认为：一个行为只要导致了加重结果，就一定构成结果加重犯。只有刑法明文规定，对加重结果加重处罚，才能构成结果加重犯。

[总结]\CJKunderwave{易考的结果加重犯}：

第一类：对加重结果持过失心理

1.故意伤害罪致人死亡。（第234条）

2.非法拘禁罪致人重伤、死亡。（第238条）

3.暴力干涉婚姻自由罪致人死亡。（第257条）

4.虐待罪致人死亡。（第260条）。\CJKunderwave{注意}，刑法在侮辱罪、诽谤罪、遗弃罪中均没有规定“若致人死亡要加重处罚”，因此这些罪中不存在结果加重犯。（2002年第43题）

5.抢夺罪（过失）致人重伤、死亡。这是司法解释的规定。

第二类：对加重结果既可以持过失心理，也可以持故意心理。

1.抢劫罪致人重伤、死亡。（第263条）

2.强奸罪致人重伤、死亡（第236条）。注意，刑法在强制猥亵罪中没有规定“致人重伤、死亡要加重处罚”，因此强制猥亵罪致人重伤、死亡不是结果加重犯。（2015年第8题）

% Source PDF physical page 56; see audit/source_coverage.json.
3.拐卖妇女、儿童罪致人重伤、死亡。（第240条）

4.放火罪致人重伤、死亡。（第115条）

例1，甲为抢劫乙，将乙杀死，拿走财物，构成抢劫罪（故意）致人死亡（2017年第2题）。

例2，甲欲强奸妇女，使用暴力不顾其死活地将妇女打成重伤昏迷，进而奸淫，之后妇女因为伤势过重死亡，甲构成强奸罪（故意）致人死亡。

\subsection*{二、一个行为}

结果加重犯是一个行为，而不是两个行为。

[注意］结果加重犯与想象竞合犯的关系。

二者相同点：行为结构相同，均是一个行为同时触犯两个罪名（A罪与B罪）。

二者区别：有无法律特别规定。想象竞合犯无法律特别规定，正常处理即可，也即在A罪与B罪中择一重罪论处；而结果加重犯是法律特别规定的产物，也即定基本犯（A罪），将B罪作为A罪的法定刑升格条件。也即，想象竞合犯是原生态的做法，结果加重犯是法律特别规定的特殊产物。

例如，甲用绳子捆绑了人质乙，出去向乙的家人要钱，回来发现乙死亡，原来绳子捆得太紧，导致乙室息死亡。甲的一个绑架行为同时触犯绑架罪（A罪）和过失致人死亡罪（B罪）。按照《刑法修正案（九）》之前的法条规定，过失致人死亡罪作为绑架罪的法定刑升格条件，加重处罚到死刑，生成结果加重犯“绑架罪（过失）致人死亡”。2015年《刑法修正案（九）》删除了该规定，那么该结果加重犯便要还原为想象竞合犯，也即甲的一个绑架行为同时触犯绑架罪和过失致人死亡罪，想象竞合，择一重罪论处，定普通的绑架罪。

\InsertDiagram{figures/manual_p056_1.tex}

\subsection*{三、因果关系}

这是指，基本犯的实行行为与加重结果之间必须具有因果关系。

\subsubsection*{（一）“因”的辨认}

\CJKunderwave{加重结果必须是基本犯的实行行为导致的，而不能是其他犯罪行为导致的}。加重结果一般都是“致人重伤、死亡”。能够致人重伤、死亡的行为大多是暴力行为。该暴力行为只有是基本犯的实行行为，其制造的加重结果，才能成为结果加重犯。该暴力行为是不是基本犯的实行行为，可根据“\CJKunderwave{行为与目的同时存在原则}”来判断。带着基本犯的目的实施的暴力行为才是基本犯的实行行为，带着其他犯罪目的实施的暴力行为是其他犯罪的实行行为。

例1，甲抢劫到乙的财物后，又为了泄愤，将乙打成重伤。甲将乙打成重伤时，没有带着取财目的，因此不属于抢劫罪的实行行为，而是故意伤害罪的实行行为。甲不构成抢劫罪（故意）致人重伤，而构成抢劫罪既遂和故意伤害罪既遂，并罚。

例2，甲欲抢劫乙，乙反抗，甲为了抢劫到财物，将乙打成重伤，压制反抗，拿走财物。甲将乙打成重伤时，带着取财目的，所以属于抢劫罪的实行行为，构成抢劫罪（故意）致人重伤。

% Source PDF physical page 57; see audit/source_coverage.json.
[巩固练习]题1，甲欲强奸乙女，乙反抗，甲为了制服乙，将乙打成重伤，进而奸淫。甲的伤害行为是为了强奸，属于强奸罪的实行行为，构成强奸罪（故意）致人重伤。

题2，甲强奸了乙女，乙女怒骂甲，甲生气，将乙打成重伤。甲不构成强奸罪致人重伤，而构成强奸罪与故意伤害罪，数罪并罚。

题3（2007年第12题），甲欲强奸乙女，乙激烈反抗，咬断甲的舌头，甲极其愤怒，将乙直接掐死。甲的杀人行为不是为了强奸妇女，而是为了泄愤报复，不是强奸罪的实行行为，而是故意杀人罪的实行行为。因此，甲不构成强奸罪（故意）致人死亡，而构成强奸罪未遂和故意杀人罪既遂，并罚。注意，甲若带着奸尸的目的而杀死乙女，该杀人行为不是强奸罪的实行行为，因为强奸罪实行行为的对象是活着的妇女。

题4（2023年试题），甲欲强奸乙女，练过散打的乙将甲制服，欲将甲扭送至公安机关。甲为逃跑，掏出弹簧刀将乙捅成重伤。甲构成强奸罪未遂和故意伤害罪既遂，并罚。

\subsubsection*{（二）因果关系的判断}

考试往往考查存在介入因素的案件，对此需要应用“介入因素两步走”标准来判断。详见后文第七讲“因果关系”。

例如，甲重伤乙后，感到后悔，送乙去医院，途中因丙违章导致车祸，车祸导致乙死亡。由于甲的伤害行为与乙的死亡无因果关系，甲仅构成故意伤害罪既遂，不构成故意伤害罪致人死亡这一结果加重犯。死亡结果应归属于丙的违章行为。
